% CUMCM 全国大学生数学建模竞赛论文模板
% 使用方法: xelatex paper.tex

\documentclass[12pt,a4paper]{article}

% ---- 中文支持 ----
\usepackage[UTF8]{ctex}
\setCJKmainfont{SimSun}[AutoFakeBold=2.5]
\setCJKsansfont{SimHei}
\setCJKmonofont{FangSong}

% ---- 数学 ----
\usepackage{amsmath,amssymb,amsthm}
\newtheorem{theorem}{定理}
\newtheorem{definition}{定义}

% ---- 图表 ----
\usepackage{graphicx}
\usepackage{float}
\usepackage{booktabs}
\usepackage{longtable}
\usepackage{array}
\usepackage{multirow}
\usepackage{subcaption}

% ---- 排版 ----
\usepackage{geometry}
\geometry{a4paper, margin=2.5cm}
\usepackage{hyperref}
\usepackage{enumitem}
\usepackage{fancyhdr}
\usepackage{lastpage}

% ---- 标题格式 ----
\usepackage{titlesec}
\titleformat{\section}{\centering\LARGE\bfseries}{}{0em}{}
\titleformat{\subsection}{\Large\bfseries}{}{0em}{}

% ---- 页眉页脚 ----
\pagestyle{fancy}
\fancyhf{}
\fancyhead[C]{数学建模竞赛论文}
\fancyfoot[C]{\thepage}

\begin{document}

% ====== 标题 ======
\title{\textbf{\Huge 数学建模竞赛论文}}
\author{}
\date{}
\maketitle

% ====== 摘要 ======
\begin{abstract}
\noindent
本文针对\ldots\ldots 问题，建立了\ldots\ldots 模型。

% 关键词
\vspace{1em}
\noindent\textbf{关键词：}数学模型；优化；评价；灵敏度分析
\end{abstract}

% ====== 问题重述 ======
\section{问题重述}

% ====== 问题分析 ======
\section{问题分析}

% ====== 模型假设与符号说明 ======
\section{模型假设与符号说明}
\subsection{模型假设}
\subsection{符号说明}

% ====== 模型建立与求解 ======
\section{模型建立与求解}

% ====== 灵敏度分析 ======
\section{灵敏度分析}

% ====== 模型评价与推广 ======
\section{模型评价与推广}
\subsection{模型优点}
\subsection{模型不足}
\subsection{模型推广}

% ====== 参考文献 ======
\begin{thebibliography}{99}
\bibitem{ref1} 姜启源, 谢金星, 叶俊. 数学模型（第五版）. 高等教育出版社, 2018.
\bibitem{ref2} 司守奎, 孙玺菁. 数学建模算法与应用（第三版）. 国防工业出版社, 2021.
\end{thebibliography}

% ====== 附录 ======
\section*{附录}

\end{document}
